\documentclass[12pt]{article}
\usepackage[utf8]{inputenc}
\usepackage[T1]{fontenc}
\usepackage[margin=1in]{geometry}
\usepackage{setspace}
\usepackage{xcolor}
\usepackage[colorlinks=true, urlcolor=blue!60!black, linkcolor=blue!60!black]{hyperref}
\setlength{\parindent}{0em}
\setlength{\parskip}{0.6em}
\renewcommand{\baselinestretch}{1.0}
\begin{document}
\begin{center}
{\LARGE \textit{Hicklin v. Orbeck}, 437 U.S.\ 518 (1978)\par}
\vspace{0.5em}
{\large Case Brief --- Two-Page Edition\par}
\end{center}

\section{Memory Jogger}
``Alaska Hire'' violated Article IV's Privileges and Immunities Clause: owning oil and gas did not let Alaska force private employers to prefer residents (\href{https://supreme.justia.com/cases/federal/us/437/518/}{437 U.S.\ 518}).

\section{Detailed Case Facts}
Citing ``uniquely high unemployment,'' Alaska required all state oil and gas leases and pipeline permits to mandate hiring ``qualified Alaska residents'' first, reaching contractors, subcontractors, suppliers, and refineries. When enforcement began at the peak of Trans-Alaska Pipeline construction, Sidney Hicklin and other workers without resident cards were barred from pipeline jobs.

\section{Procedural History}
The Alaska Superior Court upheld the Act. The Alaska Supreme Court struck the one-year residency requirement but, 3--2, upheld the resident preference based on state resource ownership. The U.S. Supreme Court reversed.

\section{Judicial Votes}
Reversed, 9--0, in an opinion by Justice Brennan joined by all Justices; no separate opinions.

\section{Holding}
A State may discriminate against nonresidents in private employment only if they are a ``peculiar source of the evil'' and the discrimination is substantially related to it. State ownership of a resource is a factor, not an exemption, and cannot justify pervasive discrimination by attenuated private employers.

\section{Analysis of Opinions}
\textbf{Majority:} Under \textit{Toomer v. Witsell}, Alaska's unemployment stemmed from lack of education and training and remoteness, not nonresidents, and the flat preference helped skilled, employed residents as much as the unemployed. Under \textit{Baldwin}, ownership is one factor, and the Act reached employers with no tie to state land or funds. Oil bound for interstate commerce reinforced the result, since the States ``must sink or swim together.''

\section{Examples: Future Applications}
\textbf{Same-Side:} a Wyoming wind-farm law requiring lessees' suppliers to hire residents first is invalid. \textbf{Opposite-Side:} higher recreational hunting fees for nonresidents are permissible under \textit{Baldwin}. \textbf{Fence-Sitter:} a city's 40\% resident quota on city-funded public works, as in \textit{Camden}.

\section{Critique}
Some originalists fault the Court for imposing ``unfettered national markets for occupations'' on the States, while others question its atextual ``fundamental rights'' filter. Doctrinal critics note that blending Article IV with Commerce Clause cases left uncertainty, and that States have few tools to aid undertrained local workers.

\section{Key Quotations}
\begin{quote}
``In sum, the Act is an attempt to force virtually all businesses that benefit in some way from the economic ripple effect of Alaska's decision to develop its oil and gas resources to bias their employment practices in favor of the State's residents. We believe that Alaska's ownership of the oil and gas that is the subject matter of Alaska Hire simply constitutes insufficient justification for the pervasive discrimination against nonresidents that the Act mandates.''
\end{quote}
\noindent --- Justice Brennan, Opinion of the Court, 437 U.S.\ at 531.

\end{document}
